\caption{\textbf{Kindness and the reference-dependent draw payoff near zero.}
The horizontal axis is the chooser's actual draw payoff,
$r_G=K(1/2-E)$, in rating points. The sample comprises competitive positions
(chooser-oriented board evaluation at least $-100$ centipawns) with $|r_G|<1$,
using the established complete-case and non-singleton chooser-by-pool sample.
The opportunity-level linear probability model includes chooser-by-speed-pool
fixed effects, continuous piecewise-linear splines in the displayed rating gap
and expected-score imbalance (with no knot at zero), a spline in the log
forgone-win premium, game-stage bands, both players' rating deviations, board
evaluation, chooser rating, both clocks and missing-clock indicators, month,
and tournament status.
Circles are means of adjusted observations in 0.25-point payoff bins,
positioned at the within-bin mean of actual $r_G$; the diamond denotes exact
zero payoff. Each adjusted observation equals the residual from the local
model plus its fitted payoff component, centered and shifted so the sample
mean equals the observed local kind-draw rate. Lines are the separate loss-
and gain-side payoff components from that same model and normalization.
The model allows a level difference between sides and a separate exact-zero
level; neither line is connected to the zero-payoff point. Bars are pointwise
95\% chooser-clustered CR1 intervals that account for estimated controls and
fixed effects, conditional on the finite model and the sample-mean display
normalization. These binned partial residuals visualize the specified local
model; they do not independently estimate an unrestricted dose-response
function. Estimates are conditional associations, not causal elasticities.}
